% ============================================================
% Linear Algebra Tutoring — shared PDF design system
% Boxed definitions/facts/hints/cautions, example cards, ruled
% working space. Loaded via pandoc -H for every week's PDF.
% ============================================================
\usepackage{fontspec}
\setmainfont{DejaVu Serif}
\usepackage[most]{tcolorbox}
\usepackage{enumitem}
\usepackage{amssymb}
\usepackage{xcolor}
\usepackage{adjustbox}
\usepackage{array}
\usepackage{longtable}
\usepackage{booktabs}
% Fragment-only pandoc->latex conversion (no --standalone) doesn't define its
% own template helper macros, needed for any table pandoc renders inside a box.
\providecommand{\real}[1]{#1}
\providecommand{\tightlist}{\setlength{\itemsep}{0pt}\setlength{\parskip}{0pt}}
% \linewidth is not reliably re-scoped inside a breakable tcolorbox (it can
% still report the page's full text width), so define a fixed, conservative
% width for auto-shrinking long display equations instead of trusting it.
\newlength{\boxmathwidth}
\AtBeginDocument{%
  \setlength{\boxmathwidth}{\textwidth}
  \addtolength{\boxmathwidth}{-30pt}%
}

% ---- palette ----------------------------------------------
\definecolor{defInk}{HTML}{1F3B73}
\definecolor{defBg}{HTML}{EEF2FB}
\definecolor{factInk}{HTML}{0F6E64}
\definecolor{factBg}{HTML}{E9F7F5}
\definecolor{hintInk}{HTML}{93650A}
\definecolor{hintBg}{HTML}{FDF3E0}
\definecolor{cautionInk}{HTML}{9C2B23}
\definecolor{cautionBg}{HTML}{FBEAE8}
\definecolor{exampleInk}{HTML}{3A3A3A}
\definecolor{exampleBg}{HTML}{F4F4F4}
\definecolor{noteInk}{HTML}{4A4A4A}
\definecolor{noteBg}{HTML}{F0F0F0}
\definecolor{ruleGray}{HTML}{BBBBBB}
\definecolor{titleRule}{HTML}{1F3B73}

% ---- boxes ---------------------------------------------------
\newtcolorbox{definitionbox}[1]{
  enhanced, breakable, colback=defBg, colframe=defInk,
  coltitle=white, colbacktitle=defInk, fonttitle=\bfseries,
  title={\footnotesize DEFINITION \quad\normalsize #1},
  boxrule=0.4pt, arc=2pt, left=5pt, right=5pt, top=5pt, bottom=7pt,
  before skip=10pt plus 2pt, after skip=10pt plus 2pt,
  fontupper=\small,
}
\newtcolorbox{factbox}[1]{
  enhanced, breakable, colback=factBg, colframe=factInk,
  coltitle=white, colbacktitle=factInk, fonttitle=\bfseries,
  title={\footnotesize KEY FACT \quad\normalsize #1},
  boxrule=0.4pt, arc=2pt, left=5pt, right=5pt, top=5pt, bottom=7pt,
  before skip=10pt plus 2pt, after skip=10pt plus 2pt,
  fontupper=\small,
}
\newtcolorbox{methodbox}[1]{
  enhanced, breakable, colback=factBg, colframe=factInk,
  coltitle=white, colbacktitle=factInk, fonttitle=\bfseries,
  title={\footnotesize METHOD \quad\normalsize #1},
  boxrule=0.4pt, arc=2pt, left=5pt, right=5pt, top=5pt, bottom=7pt,
  before skip=10pt plus 2pt, after skip=10pt plus 2pt,
  fontupper=\small,
}
\newtcolorbox{notebox}[1]{
  enhanced, breakable, colback=noteBg, colframe=noteInk!60,
  coltitle=white, colbacktitle=noteInk, fonttitle=\bfseries,
  title={\footnotesize NOTE \quad\normalsize #1},
  boxrule=0.4pt, arc=2pt, left=5pt, right=5pt, top=5pt, bottom=7pt,
  before skip=10pt plus 2pt, after skip=10pt plus 2pt,
  fontupper=\small,
}
\newtcolorbox{examplebox}[1]{
  enhanced, breakable, colback=exampleBg, colframe=exampleInk!70,
  coltitle=white, colbacktitle=exampleInk, fonttitle=\bfseries,
  title={\footnotesize EXAMPLE~#1},
  boxrule=0.3pt, arc=2pt, left=6pt, right=6pt, top=6pt, bottom=8pt,
  before skip=12pt plus 2pt, after skip=12pt plus 2pt,
  fontupper=\small,
}
\newtcolorbox{hintbox}{
  enhanced, breakable, colback=hintBg, colframe=hintInk,
  boxrule=0pt, leftrule=3pt, arc=1pt, left=9pt, right=9pt, top=5pt, bottom=5pt,
  before skip=8pt plus 2pt, after skip=4pt,
}
\newtcolorbox{cautionbox}{
  enhanced, breakable, colback=cautionBg, colframe=cautionInk,
  boxrule=0pt, leftrule=3pt, arc=1pt, left=9pt, right=9pt, top=5pt, bottom=5pt,
  before skip=4pt, after skip=10pt plus 2pt,
}
\newcommand{\hintlabel}{\textbf{\textcolor{hintInk}{HINT}}\quad}
\newcommand{\cautionlabel}{\textbf{\textcolor{cautionInk}{CAUTION --- before you solve}}\quad}

% ---- problem number badge + working-space rules ---------------
\newtcolorbox{problembadge}[1]{
  enhanced, colback=defInk, colframe=defInk, coltext=white,
  boxrule=0pt, arc=10pt, left=0pt,right=0pt,top=0pt,bottom=0pt,
  width=fit, height=2.1em, valign=center, halign=center,
  fontupper=\bfseries\large,
}
\newcommand{\worklines}[1]{%
  \vspace{2pt}\par\noindent{\small\itshape\color{gray}Your working:}\\[4pt]
  \foreach \n in {1,...,#1}{\textcolor{ruleGray}{\hrulefill}\\[0.62cm]}
  \vspace{2pt}
}
\usepackage{pgffor}

% ---- title block -----------------------------------------------
\newcommand{\worksheettitle}[3]{%
  {\centering
   {\fontsize{28}{32}\selectfont\bfseries Linear Algebra --- Week #1}\par
   \vspace{4pt}
   {\color{titleRule}\rule{\linewidth}{1.4pt}}\par
   \vspace{6pt}
   {\large\itshape #2}\par
   \vspace{2pt}
   {\small\color{gray} #3}\par
  }
  \vspace{14pt}
}

\newcommand{\sectionbanner}[1]{%
  \vspace{10pt}\noindent{\Large\bfseries\color{defInk} #1}\par
  \noindent{\color{defInk}\rule{\linewidth}{1pt}}\vspace{10pt}%
}

% ---- misc: tighter lists, no section numbers --------------------
\setlist[enumerate]{leftmargin=1.6em}
\setcounter{secnumdepth}{0}
\pagestyle{plain}
